% Generated draft: requires mathematical and source review.
\documentclass[11pt]{article}
\usepackage[margin=25mm]{geometry}
\usepackage{fontspec}
\setmainfont{Latin Modern Roman}
\usepackage{amsmath,amssymb,mathtools,unicode-math}
\setmathfont{Latin Modern Math}
\usepackage{enumitem,needspace,xurl,hyperref}
\setlength{\parindent}{0pt}
\setlength{\parskip}{5pt}
\setlength{\emergencystretch}{3em}
\newcommand{\N}{\mathbb{N}}
\newcommand{\Z}{\mathbb{Z}}
\newcommand{\Q}{\mathbb{Q}}
\newcommand{\R}{\mathbb{R}}
\newcommand{\C}{\mathbb{C}}
\newcommand{\F}{\mathbb{F}}
\newcommand{\A}{\mathbb{A}}
\newcommand{\PP}{\mathbb{P}}
\newcommand{\E}{\mathbb{E}}
\newcommand{\eps}{\varepsilon}
\newcommand{\dd}{\,\mathrm{d}}
\begin{document}
\section*{UnsolvedMath ID: 700799}
\textbf{Source problem code:} catalogue entry
\subsection{Definitions and mathematical statement}
A group \(G\) is locally finite if every subgroup generated by finitely many elements of \(G\) is finite. It is simple if \(G\neq\{1\}\) and its only normal subgroups are \(\{1\}\) and \(G\). A group \(H\) is residually finite if, for every \(h\in H\setminus\{1\}\), there exist a finite group \(F\) and a homomorphism \(\varphi:H\to F\) such that \(\varphi(h)\neq 1\).

Using the usual unrestricted-field convention, a group \(G\) is linear if there exist a field \(K\), an integer \(n\geq 1\), and an injective homomorphism \(G\to\operatorname{GL}_n(K)\). It is non-linear if no such field, integer, and homomorphism exist. The supplied record does not explicitly specify this convention or restrict the coefficient field; this interpretation requires confirmation against the source.

\textbf{Problem (R. Phillips, 1990).} Does there exist a non-linear, locally finite, simple group \(G\) such that every proper subgroup \(H<G\) is residually finite? Explicitly, the requirement on proper subgroups is
\[
\forall H<G\;\forall h\in H\setminus\{1\}\;\exists F\text{ finite}\;\exists\varphi:H\to F\text{ a homomorphism},\qquad \varphi(h)\neq 1.
\]
\subsection{Short English statement}
Does there exist a simple group in which every finitely generated subgroup is finite, which has no faithful finite-dimensional representation over any field, and in which every nonidentity element of every proper subgroup survives in some finite quotient of that subgroup? The interpretation of non-linearity uses the convention stated above.
\subsection{Sources}
\begin{itemize}
\item Supplied local UnsolvedMath registry record: permanent ID \(700799\), exact source code \texttt{catalogue entry}, title \emph{Kourovka Notebook Problem 11.107}. The supplied status is \texttt{open}; this is a registry label, not an independently verified mathematical assessment.
\item \emph{The Kourovka Notebook (Unsolved Problems in Group Theory)}, Problem 11.107, attributed to R. Phillips, 11th issue (1990). The supplied record identifies the 21st issue edition (2026), \texttt{arXiv:1401.0300v48}.
\item The supplied collection description names E. I. Khukhro and V. D. Mazurov as editors of the Kourovka Notebook collection covering Issues 1--21 (1965--2026). These bibliographic details are supplied attribution and have not been independently checked.
\end{itemize}
\end{document}
